\bta{2024}

\begin{enumerate}
	%\renewcommand{\labelenumi}{\arabic{enumi}.}
	% A(\Alph) a(\alph) I(\Roman) i(\roman) 1(\arabic)
	%设定全局标号series=example	%引用全局变量resume=example
	%[topsep=-0.3em,parsep=-0.3em,itemsep=-0.3em,partopsep=-0.3em]
	%可使用leftmargin调整列表环境左边的空白长度 [leftmargin=0em]
	\item
选择题 (本题满分 50 分，每小题 5 分)
\begin{enumerate}
	%\renewcommand{\labelenumi}{\arabic{enumi}.}
	% A(\Alph) a(\alph) I(\Roman) i(\roman) 1(\arabic)
	%设定全局标号series=example	%引用全局变量resume=example
	%[topsep=-0.3em,parsep=-0.3em,itemsep=-0.3em,partopsep=-0.3em]
	%可使用leftmargin调整列表环境左边的空白长度 [leftmargin=0em]
	\item
下列极限不存在或为无穷大的是 \xzanswer{B}

\fourchoices
{$\lim\limits_{x \rightarrow 0} x \sin \frac{1}{x^2}$}
{$\lim\limits_{x \rightarrow+\infty}\left(1+\frac{\ln x}{x}\right)^x$}
{$\lim\limits_{x \rightarrow 0} e^{\frac{\sin x^2}{x}}$}
{$\lim\limits_{x \rightarrow+\infty} \frac{x^2}{e^x+\arctan x}$}


\item
设 $f(x)=\int_0^x(t-2) e^{t^2} d t$，则 $f(x)$ 在 $x \in(1,4)$ 中 \xzanswer{A}

\fourchoices
{有极小值点}
{有极大值点}
{单调递增}
{单调递减}


\item
已知向量 $|\vec{a}|=1$，$|\vec{b}|=2$，$\vec{a} \cdot \vec{b}=\sqrt{3}$，
则以 $\vec{a}$ 和 $\vec{b}$ 为邻边的平行四边形的面积为 \xzanswer{C}

\fourchoices
{$\sqrt{3}$}
{$\sqrt{2}$}
{$1$}
{$\frac{1}{2}$}


\item
设函数 $f(x)$ 可导，$F(x)=f(x)(1+\sin |x|)$，
则 $f(0)=0$ 是 $F(x)$ 在 $x=0$ 处可导的 \xzanswer{A}

\fourchoices
{充分必要条件}
{充分但非必要条件}
{必要但非充分条件}
{既非充分条件又非必要条件}


\item
以下四个命题中，正确的是 \xzanswer{C}

\fourchoices
{若 $f^{\prime}(x)$ 在 $(0,1)$ 内连续，则 $f(x)$ 在 $(0,1)$ 内有界}
{若 $f(x)$ 在 $(0,1)$ 内连续，则 $f(x)$ 在 $(0,1)$ 内有界}
{若 $f^{\prime}(x)$ 在 $(0,1)$ 内有界，则 $f(x)$ 在 $(0,1)$ 内有界}
{若 $f(x)$ 在 $(0,1)$ 内有界且可导，则 $f^{\prime}(x)$ 在 $(0,1)$ 内有界}


\item
常微分方程 $\left(x-2 x y-y^2\right) \frac{d y}{d x}+y^2=0$ 的通解为 \xzanswer{D}

\fourchoices
{$y=x^2+C x^2 e^{\frac{1}{y}}$}
{$y=C x^2 e^{\frac{1}{y}}$}
{$y=x^2+C x^2 e^{\frac{3}{y}}$}
{$x=y^2+C y^2 e^{\frac{1}{y}}$}


\item
设函数 $f(t)$ 二次连续可微，若令 $t=x y$，则 $\frac{\partial^2 f}{\partial x \partial y}$ 的表达式是 \xzanswer{C}

\fourchoices
{$t f^{\prime \prime}(t)-f^{\prime}(t)$}
{$f^{\prime \prime}(t)-t f^{\prime}(t)$}
{$t f^{\prime \prime}(t)+f^{\prime}(t)$}
{$f^{\prime \prime}(t)+t f^{\prime}(t)$}


\item
设 $a_n \neq 0\ (n=1,2, \ldots)$，且 $\lim\limits_{n \rightarrow \infty} a_n=a \neq 0$，
则下列说法正确的是 \xzanswer{C}

\fourchoices
{级数 $\sum\limits_{n=1}^{\infty}\left|a_{n+1}-a_n\right|$ 收敛，
而级数 $\sum\limits_{n=1}^{\infty}\left|\frac{1}{a_{n+1}}-\frac{1}{a_n}\right|$ 发散}
{级数 $\sum\limits_{n=1}^{\infty}\left|\frac{1}{a_{n+1}}-\frac{1}{a_n}\right|$ 收敛，
而级数 $\sum\limits_{n=1}^{\infty}\left|a_{n+1}-a_n\right|$ 发散}
{级数 $\sum\limits_{n=1}^{\infty}\left|a_{n+1}-a_n\right|$ 和
$\sum\limits_{n=1}^{\infty}\left|\frac{1}{a_{n+1}}-\frac{1}{a_n}\right|$ 具有相同的敛散性}
{级数 $\sum\limits_{n=1}^{\infty}\left|a_{n+1}-a_n\right|$ 和
$\sum\limits_{n=1}^{\infty}\left|\frac{1}{a_{n+1}}-\frac{1}{a_n}\right|$ 的敛散性没有必然联系}


\item
设 $f(x)$ 是实数轴上的连续函数，下列说法正确的是 \xzanswer{B}

\fourchoices
{若 $f(x)$ 是奇函数，则它的原函数一定为奇函数}
{若 $f(x)$ 是奇函数，则它的原函数一定为偶函数}
{若 $f(x)$ 是偶函数，则它的原函数一定为奇函数}
{若 $f(x)$ 是周期函数，则它的原函数一定为周期函数}


\item
下列广义积分发散的是 \xzanswer{D}

\fourchoices
{$\int_0^1 \frac{\ln x}{1-x} d x$}
{$\int_0^1 \frac{\cos x^2 \cdot \ln x}{x} d x$}
{$\int_0^1 \frac{\ln x}{1-x^2} d x$}
{$\int_0^1 \frac{\ln x}{(1-x)^2} d x$}


\end{enumerate}


\banswer{
\begin{enumerate}
	%\renewcommand{\labelenumi}{\arabic{enumi}.}
	% A(\Alph) a(\alph) I(\Roman) i(\roman) 1(\arabic)
	%设定全局标号series=example	%引用全局变量resume=example
	%[topsep=-0.3em,parsep=-0.3em,itemsep=-0.3em,partopsep=-0.3em]
	%可使用leftmargin调整列表环境左边的空白长度 [leftmargin=0em]
	\item
	\textbf{B}。A 中 $\left|x \sin \frac{1}{x^2}\right| \leqslant|x| \rightarrow 0$，极限为 $0$；
	C 中 $\frac{\sin x^2}{x} \sim \frac{x^2}{x}=x \rightarrow 0$，极限为 $e^0=1$；
	D 中分母以 $e^x$ 为主，极限为 $0$。而 B 中取对数
	\[
	x \ln \left(1+\frac{\ln x}{x}\right)=x\left(\frac{\ln x}{x}-\frac{\ln ^2 x}{2 x^2}+o\left(\frac{\ln ^2 x}{x^2}\right)\right)
	=\ln x-\frac{\ln ^2 x}{2 x}+o(1) \rightarrow+\infty,
	\]
	故 B 中极限为 $+\infty$。

	\item
	\textbf{A}。$f^{\prime}(x)=(x-2) e^{x^2}$，在 $(1,2)$ 上 $f^{\prime}<0$，在 $(2,4)$ 上 $f^{\prime}>0$，
	故 $x=2$ 为 $(1,4)$ 内的极小值点，且 $f$ 在 $(1,4)$ 上不单调。

	\item
	\textbf{C}。由 $\cos \theta=\frac{\vec{a} \cdot \vec{b}}{|\vec{a}||\vec{b}|}=\frac{\sqrt{3}}{2}$ 得 $\theta=\frac{\pi}{6}$，
	于是 $S=|\vec{a}||\vec{b}| \sin \theta=1 \times 2 \times \frac{1}{2}=1$。

	\item
	\textbf{A}。$F(x)=f(x)+f(x) \sin |x|$，其中 $f$ 处处可导，故只需考察 $g(x)=f(x) \sin |x|$。
	$g(0)=0$，且
	\begin{align*}
	\lim _{x \rightarrow 0^{+}} \frac{g(x)-g(0)}{x}&=\lim _{x \rightarrow 0^{+}} f(x) \frac{\sin x}{x}=f(0),\\
	\lim _{x \rightarrow 0^{-}} \frac{g(x)-g(0)}{x}&=\lim _{x \rightarrow 0^{-}} f(x) \frac{-\sin x}{x}=-f(0).
	\end{align*}
	故 $g^{\prime}(0)$ 存在 $\Longleftrightarrow f(0)=-f(0) \Longleftrightarrow f(0)=0$，即为充分必要条件。

	\item
	\textbf{C}。A、B 的反例同为 $f(x)=\frac{1}{x}$：它与 $f^{\prime}(x)=-\frac{1}{x^2}$ 都在 $(0,1)$ 内连续，
	但 $f$ 无界。D 的反例为 $f(x)=\sin \frac{1}{x}$：$f$ 有界可导，
	而 $f^{\prime}(x)=-\frac{1}{x^2} \cos \frac{1}{x}$ 无界。
	C 正确：若 $\left|f^{\prime}\right| \leqslant M$ 于 $(0,1)$，由 Lagrange 中值定理，
	$\forall x \in(0,1)$ 有 $\left|f(x)-f\left(\frac{1}{2}\right)\right|=\left|f^{\prime}(\xi)\right|\left|x-\frac{1}{2}\right| \leqslant \frac{M}{2}$，
	故 $|f(x)| \leqslant\left|f\left(\frac{1}{2}\right)\right|+\frac{M}{2}$ 有界。

	\item
	\textbf{D}。原方程即 $y^2 d x+\left(x(1-2 y)-y^2\right) d y=0$，把 $x$ 视为 $y$ 的函数，得一阶线性方程
	\[
	\frac{d x}{d y}+\frac{1-2 y}{y^2} x=1 .
	\]
	因 $\int\left(\frac{1}{y^2}-\frac{2}{y}\right) d y=-\frac{1}{y}-2 \ln |y|$，取积分因子 $\mu=y^{-2} e^{-1 / y}$，则
	$(\mu x)^{\prime}=y^{-2} e^{-1 / y}$。由 $u=-\frac{1}{y}$ 得 $\int y^{-2} e^{-1 / y} d y=e^{-1 / y}$，
	故 $y^{-2} e^{-1 / y} x=e^{-1 / y}+C$，即 $x=y^2+C y^2 e^{\frac{1}{y}}$。

	\item
	\textbf{C}。$\frac{\partial f}{\partial x}=y f^{\prime}(t)$，故
	\[
	\frac{\partial^2 f}{\partial x \partial y}=\frac{\partial}{\partial y}\left[y f^{\prime}(x y)\right]
	=f^{\prime}(t)+x y f^{\prime \prime}(t)=t f^{\prime \prime}(t)+f^{\prime}(t) .
	\]

	\item
	\textbf{C}。由
	$\left|\frac{1}{a_{n+1}}-\frac{1}{a_n}\right|=\frac{\left|a_{n+1}-a_n\right|}{\left|a_n a_{n+1}\right|}$
	及 $a_n a_{n+1} \rightarrow a^2 \neq 0$，得
	\[
	\lim _{n \rightarrow \infty} \frac{\left|\frac{1}{a_{n+1}}-\frac{1}{a_n}\right|}{\left|a_{n+1}-a_n\right|}=\frac{1}{a^2} \in(0,+\infty).
	\]
	由正项级数的比较判别法（极限形式），两级数同敛散。

	\item
	\textbf{B}。$f$ 的全部原函数为 $F(x)=\int_0^x f(t) d t+C$。若 $f$ 为奇函数，
	则 $\int_0^x f(t) d t$ 为偶函数，偶函数加常数仍为偶函数，故 B 正确；
	此时 $F$ 一般不是奇函数（若既奇又偶则 $F \equiv 0$），A 错。
	若 $f$ 为偶函数，$\int_0^x f$ 为奇函数，但 $C \neq 0$ 时 $F$ 不再是奇函数，
	故 C 中“一定”不成立。D 的反例为 $f(x) \equiv 1$，其原函数 $x+C$ 非周期函数。

	\item
	\textbf{D}。A 中 $x \rightarrow 1$ 时 $\frac{\ln x}{1-x} \rightarrow-1$，$x \rightarrow 0^{+}$ 时 $\ln x$ 可积，收敛；
	C 中 $x \rightarrow 1$ 时 $\frac{\ln x}{1-x^2} \rightarrow-\frac{1}{2}$，同样收敛。
	D 中 $x \rightarrow 1^{-}$ 时 $\ln x \sim-(1-x)$，故被积函数 $\sim-\frac{1}{1-x}$，
	而 $\int_0^1 \frac{d x}{1-x}$ 发散，故 D 发散。

	\par\vspace{1mm}\noindent
	\textbf{注：}{\kaishu
	按题面，B 其实也发散：$x \rightarrow 0^{+}$ 时 $\cos x^2 \rightarrow 1$，被积函数 $\sim \frac{\ln x}{x}$，
	而 $\int_\varepsilon^1 \frac{\ln x}{x} d x=-\frac{\ln ^2 \varepsilon}{2} \rightarrow-\infty$。
	因此本题印制的四个选项中 B、D 均发散，试题本身有瑕疵；
	若按 A、C、D 同族（考察 $x \rightarrow 1$ 处的奇性）的命题意图，应选 D。
	}\par
\end{enumerate}
}


\item
(本题满分 10 分)
数列 $\left\{a_n\right\}_{n=0}^{\infty}$ 满足条件
$a_0=0$，$a_1=1$，$a_{n+1}=\frac{a_n+n a_{n-1}}{n+1}\ (n \geqslant 1)$，
求数列 $\left\{a_n\right\}_{n=0}^{\infty}$ 的极限。
(提示：试考虑 $a_{n+1}-a_n$ 的递归公式)

\banswer{
	记 $d_n=a_n-a_{n-1}\ (n \geqslant 1)$，则 $d_1=a_1-a_0=1$，且
	\[
	a_{n+1}-a_n=\frac{a_n+n a_{n-1}}{n+1}-a_n=\frac{n a_{n-1}-n a_n}{n+1}=-\frac{n}{n+1}\left(a_n-a_{n-1}\right),
	\]
	即 $d_{n+1}=-\frac{n}{n+1} d_n$。逐项递推得
	\[
	d_n=\left(-\frac{n-1}{n}\right)\left(-\frac{n-2}{n-1}\right) \cdots\left(-\frac{1}{2}\right) d_1
	=\frac{(-1)^{n-1}}{n},\qquad n \geqslant 1 .
	\]
	于是
	\[
	a_n=a_0+\sum_{k=1}^n d_k=\sum_{k=1}^n \frac{(-1)^{k-1}}{k}.
	\]
	由 Leibniz 交错级数判别法，$\sum\limits_{k=1}^{\infty} \frac{(-1)^{k-1}}{k}$ 收敛，
	且其和恰为 $\ln (1+x)$ 在 $x=1$ 处的值，故
	\[
	\lim _{n \rightarrow \infty} a_n=\ln 2 .
	\]
}


\item
(本题满分 10 分)
设 $f(x)=\left\{\begin{array}{cl}
\frac{\sin (x-a)}{x-1}, & x \neq 1 \\
b, & x=1
\end{array}\right.$，
求 $a$，$b$ 的值，使得 $f(x)$ 在 $x=1$ 连续且 $f^{\prime}(a)=\frac{1}{\pi}$。

\banswer{
	\textbf{连续性条件.} 要使 $\lim\limits_{x \rightarrow 1} \frac{\sin (x-a)}{x-1}$ 存在有限，
	分子必须在 $x=1$ 处为零，即 $\sin (1-a)=0$，故
	\[
	1-a=k \pi \quad(k \in \mathbb{Z}),\qquad a=1-k \pi .
	\]
	此时 $\sin (x-a)=\sin \bigl((x-1)+k \pi\bigr)=(-1)^k \sin (x-1)$，于是
	\[
	b=\lim _{x \rightarrow 1} \frac{(-1)^k \sin (x-1)}{x-1}=(-1)^k .
	\]

	\textbf{导数条件.} 记 $u=x-1$，则对 $x \neq 1$ 有 $f(x)=(-1)^k \frac{\sin u}{u}$，
	\[
	f^{\prime}(x)=(-1)^k \frac{u \cos u-\sin u}{u^2} .
	\]
	在 $x=a=1-k \pi$ 处 $u=-k \pi$（$k \neq 0$，否则 $a=1$ 而 $f^{\prime}(1)=0 \neq \frac{1}{\pi}$），
	代入 $\cos (-k \pi)=(-1)^k$，$\sin (-k \pi)=0$ 得
	\[
	f^{\prime}(a)=(-1)^k \cdot \frac{(-k \pi)(-1)^k}{k^2 \pi^2}=-\frac{1}{k \pi} .
	\]
	由 $-\frac{1}{k \pi}=\frac{1}{\pi}$ 得 $k=-1$。故
	\[
	\boxed{a=1+\pi,\qquad b=-1 .}
	\]

	\textbf{验证.} $f(x)=\frac{\sin (x-1-\pi)}{x-1}=-\frac{\sin (x-1)}{x-1}$，
	$\lim\limits_{x \rightarrow 1} f(x)=-1=b$，连续；
	且 $f^{\prime}(1+\pi)=-\frac{\pi \cos \pi-\sin \pi}{\pi^2}=\frac{1}{\pi}$，符合要求。
}


\item
(本题满分 10 分)
计算极限 $I=\lim\limits_{n \rightarrow+\infty} \int_0^1 n x e^{-n\left(x^2+x^3\right)} d x$。

\banswer{
	作代换 $x=\frac{t}{\sqrt{n}}$，则 $n x\, d x=t\, d t$，且
	$-n\left(x^2+x^3\right)=-\left(t^2+\frac{t^3}{\sqrt{n}}\right)$，于是
	\[
	I_n=\int_0^1 n x e^{-n\left(x^2+x^3\right)} d x=\int_0^{\sqrt{n}} t e^{-t^2} e^{-t^3 / \sqrt{n}}\, d t .
	\]
	对 $t \geqslant 0$ 有 $0<e^{-t^3 / \sqrt{n}} \leqslant 1$，故被积函数被可积函数 $t e^{-t^2}$ 控制，
	且逐点有 $e^{-t^3 / \sqrt{n}} \rightarrow 1$。由控制收敛定理（或分段估计：
	固定 $A>0$，在 $[0,A]$ 上一致收敛，在 $[A,+\infty)$ 上尾部 $\int_A^{\infty} t e^{-t^2} d t=\frac{e^{-A^2}}{2}$ 可任意小），
	\[
	I=\lim _{n \rightarrow+\infty} I_n=\int_0^{+\infty} t e^{-t^2} d t=\left[-\frac{1}{2} e^{-t^2}\right]_0^{+\infty}=\frac{1}{2} .
	\]
	\par\vspace{1mm}\noindent
	\textbf{注：}{\kaishu
	直观地看，$n \rightarrow \infty$ 时积分的质量集中在 $x=O\!\left(n^{-1 / 2}\right)$ 的小邻域内，
	那里 $x^3=o\left(x^2\right)$ 可以忽略，于是
	$I \approx \int_0^1 n x e^{-n x^2} d x=\frac{1-e^{-n}}{2} \rightarrow \frac{1}{2}$。
	}\par
}


\item
(本题满分 10 分)
设 $x_1 \cdot x_2>0$，证明：存在 $x_1$ 与 $x_2$ 之间的 $\xi$ 使得
\[
x_1 e^{x_2}-x_2 e^{x_1}=(1-\xi) e^{\xi}\left(x_1-x_2\right).
\]

\banswer{
	不妨设 $x_1 \neq x_2$（若 $x_1=x_2$，两端同为 $0$，结论显然）。
	由 $x_1 x_2>0$ 知 $x_1$ 与 $x_2$ 同号，故 $0$ 不在 $x_1$ 与 $x_2$ 之间，
	设该闭区间为 $[\alpha, \beta]$，则 $0 \notin[\alpha, \beta]$。

	在 $[\alpha, \beta]$ 上取
	\[
	g(x)=\frac{e^x}{x}, \qquad h(x)=\frac{1}{x},
	\]
	二者都连续且可导，且 $h^{\prime}(x)=-\frac{1}{x^2} \neq 0$，$h\left(x_1\right) \neq h\left(x_2\right)$。
	由 Cauchy 中值定理，存在 $\xi$ 介于 $x_1$ 与 $x_2$ 之间，使
	\[
	\frac{g\left(x_2\right)-g\left(x_1\right)}{h\left(x_2\right)-h\left(x_1\right)}=\frac{g^{\prime}(\xi)}{h^{\prime}(\xi)} .
	\]
	而
	\[
	g^{\prime}(x)=\frac{x e^x-e^x}{x^2}=\frac{(x-1) e^x}{x^2},\qquad
	\frac{g^{\prime}(x)}{h^{\prime}(x)}=\frac{(x-1) e^x}{x^2} \cdot\left(-x^2\right)=(1-x) e^x .
	\]
	故
	\[
	\frac{e^{x_2}}{x_2}-\frac{e^{x_1}}{x_1}=(1-\xi) e^{\xi}\left(\frac{1}{x_2}-\frac{1}{x_1}\right).
	\]
	两端同乘 $x_1 x_2\ (\neq 0)$ 即得
	\[
	x_1 e^{x_2}-x_2 e^{x_1}=(1-\xi) e^{\xi}\left(x_1-x_2\right).
	\]
}


\item
(本题满分 10 分)
计算函数 $f(x)=\int_0^{\cos x} \ln (2+2 t) d t$ 的导函数 $f^{\prime}(x)$
在开区间 $(0,2 \pi)$ 上零点个数。

\banswer{
	由复合函数求导（变限积分）
	\[
	f^{\prime}(x)=\ln (2+2 \cos x) \cdot(-\sin x)=-\sin x \cdot \ln (2+2 \cos x).
	\]
	在 $(0,2 \pi)$ 上令 $f^{\prime}(x)=0$：
	\begin{itemize}
		\item $\sin x=0$：得 $x=\pi$。此时 $\cos x=-1$，$\ln (2+2 \cos x)=\ln 0$ 无意义，
		需按定义补充验证。取 $x=\pi+\varepsilon$，则 $2+2 \cos x=\varepsilon^2+O\left(\varepsilon^4\right)$，
		$\sin x=-\sin \varepsilon \approx-\varepsilon$，故
		$f^{\prime}(x) \approx 2 \varepsilon \ln |\varepsilon| \rightarrow 0$；
		事实上由 $\int_{-1}^{-1+\delta} \ln (2+2 t) d t=\delta(\ln (2 \delta)-1)$ 及 $\delta=\frac{\varepsilon^2}{2}$ 可直接算出
		$f^{\prime}(\pi)=0$。故 $x=\pi$ 是零点。
		\item $\ln (2+2 \cos x)=0$：即 $2+2 \cos x=1$，$\cos x=-\frac{1}{2}$，
		得 $x=\frac{2 \pi}{3}$ 与 $x=\frac{4 \pi}{3}$。
	\end{itemize}
	综上，$f^{\prime}(x)$ 在 $(0,2 \pi)$ 上的零点为
	$x=\frac{2 \pi}{3},\ \pi,\ \frac{4 \pi}{3}$，共 $\boxed{3}$ 个。
}


\item
(本题满分 10 分)
求幂级数 $\sum\limits_{n=1}^{\infty}\left(1+\frac{1}{2}+\frac{1}{3}+\frac{1}{4}+\cdots+\frac{1}{n}\right) x^n$
的收敛域以及和函数 $S(x)$。
(提示求 $S(x)$ 时，试图将所考虑的级数表示为两个幂级数的乘积)

\banswer{
	记 $H_n=\sum\limits_{k=1}^n \frac{1}{k}$。

	\textbf{收敛域.} 因 $\frac{H_{n+1}}{H_n} \rightarrow 1$，由比值法收敛半径 $R=1$。
	端点处：$x=1$ 时通项 $H_n \rightarrow+\infty$，$x=-1$ 时通项 $(-1)^n H_n$ 不趋于 $0$，
	均不满足收敛的必要条件。故收敛域为 $(-1,1)$。

	\textbf{和函数.} 在 $(-1,1)$ 内，
	\[
	-\ln (1-x)=\sum_{n=1}^{\infty} \frac{x^n}{n},\qquad \frac{1}{1-x}=\sum_{m=0}^{\infty} x^m .
	\]
	作 Cauchy 乘积，$x^n$ 的系数为 $\sum\limits_{k=1}^n \frac{1}{k} \cdot 1=H_n$，即
	\[
	\left(\sum_{n=1}^{\infty} \frac{x^n}{n}\right)\left(\sum_{m=0}^{\infty} x^m\right)=\sum_{n=1}^{\infty} H_n x^n .
	\]
	故
	\[
	S(x)=-\frac{\ln (1-x)}{1-x},\qquad x \in(-1,1).
	\]
}


\item
(本题满分 10 分)
求 $\varphi(x)=e^x-\int_0^x(x-u) \varphi(u) d u$，其中 $\varphi(x)$ 为连续函数，求 $\varphi(x)$。

\banswer{
	把积分拆开：
	\[
	\varphi(x)=e^x-x \int_0^x \varphi(u) d u+\int_0^x u \varphi(u) d u .
	\]
	由 $\varphi$ 连续知右端可导，求导得
	\[
	\varphi^{\prime}(x)=e^x-\int_0^x \varphi(u) d u-x \varphi(x)+x \varphi(x)=e^x-\int_0^x \varphi(u) d u,
	\]
	再求导得
	\[
	\varphi^{\prime \prime}(x)=e^x-\varphi(x),\qquad \text{即}\quad \varphi^{\prime \prime}+\varphi=e^x .
	\]
	初值由原式与一阶导式在 $x=0$ 取值得到：$\varphi(0)=e^0=1$，$\varphi^{\prime}(0)=e^0=1$。

	齐次解为 $C_1 \cos x+C_2 \sin x$；设特解 $\varphi^*=A e^x$，代入得 $2 A=1$，$A=\frac{1}{2}$。故
	\[
	\varphi(x)=C_1 \cos x+C_2 \sin x+\frac{1}{2} e^x .
	\]
	由 $\varphi(0)=C_1+\frac{1}{2}=1$ 得 $C_1=\frac{1}{2}$；
	由 $\varphi^{\prime}(0)=C_2+\frac{1}{2}=1$ 得 $C_2=\frac{1}{2}$。所以
	\[
	\varphi(x)=\frac{1}{2}\left(e^x+\cos x+\sin x\right).
	\]
}


\item
(本题满分 10 分)
求由方程 $2 x^2+y^2+z^2+2 x y-2 x-2 y-4 z-4=0$ 所确定的函数 $z=z(x, y)$ 的极值。

\banswer{
	\textbf{隐函数求偏导.} 分别对 $x$、$y$ 求偏导：
	\[
	4 x+2 z z_x+2 y-2-4 z_x=0 \Longrightarrow z_x=\frac{1-2 x-y}{z-2},
	\]
	\[
	2 y+2 z z_y+2 x-2-4 z_y=0 \Longrightarrow z_y=\frac{1-x-y}{z-2} .
	\]
	驻点条件 $1-2 x-y=0$，$1-x-y=0$，两式相减得 $x=0$，进而 $y=1$。
	代入原方程得 $z^2-4 z-5=0$，即 $z=5$ 或 $z=-1$（此时 $z-2 \neq 0$，隐函数存在）。

	\textbf{二阶判别.} 在驻点处分子为零，故
	\[
	A=z_{x x}=\frac{-2}{z-2},\qquad B=z_{x y}=\frac{-1}{z-2},\qquad C=z_{y y}=\frac{-1}{z-2},
	\]
	\[
	A C-B^2=\frac{2}{(z-2)^2}-\frac{1}{(z-2)^2}=\frac{1}{(z-2)^2}>0 .
	\]
	当 $z=5$ 时 $A=-\frac{2}{3}<0$，故 $z(0,1)=5$ 为\textbf{极大值}；
	当 $z=-1$ 时 $A=\frac{2}{3}>0$，故 $z(0,1)=-1$ 为\textbf{极小值}。

	\par\vspace{1mm}\noindent
	\textbf{注：}{\kaishu
	配方可以一眼看出结论：
	$2 x^2+2 x y+y^2=x^2+(x+y)^2$，于是原方程等价于
	\[
	x^2+(x+y-1)^2+(z-2)^2=9,
	\]
	这是球心 $(0,1,2)$（在 $(x, x+y-1, z-2)$ 坐标下）、半径 $3$ 的球面，
	故 $z \in[-1,5]$，两端点分别在 $x=0$，$y=1$ 处取到。
	}\par
}


\item
(本题满分 10 分)
求点 $P_1(1,2,1)$ 关于直线
$L:\left\{\begin{array}{l}2 x+y+z=0 \\ x-y+z=1\end{array}\right.$
的对称点 $P_2$ 的坐标。

\banswer{
	\textbf{直线的点向式.} 方向向量
	\[
	\vec{s}=(2,1,1) \times(1,-1,1)=(2,-1,-3),\qquad |\vec{s}|^2=14 .
	\]
	取 $z=0$ 解得 $2 x+y=0$，$x-y=1$，即 $x=\frac{1}{3}$，$y=-\frac{2}{3}$，
	得 $L$ 上一点 $A\left(\frac{1}{3},-\frac{2}{3}, 0\right)$。

	\textbf{求垂足.} $\overrightarrow{A P_1}=\left(\frac{2}{3}, \frac{8}{3}, 1\right)$，
	\[
	\overrightarrow{A P_1} \cdot \vec{s}=\frac{4}{3}-\frac{8}{3}-3=-\frac{13}{3},\qquad
	t=\frac{\overrightarrow{A P_1} \cdot \vec{s}}{|\vec{s}|^2}=-\frac{13}{42} .
	\]
	于是垂足
	\[
	F=A+t \vec{s}=\left(\frac{1}{3}-\frac{26}{42},\ -\frac{2}{3}+\frac{13}{42},\ \frac{39}{42}\right)
	=\left(-\frac{2}{7},\ -\frac{5}{14},\ \frac{13}{14}\right).
	\]

	\textbf{求对称点.} $F$ 为 $P_1 P_2$ 的中点，故
	\[
	P_2=2 F-P_1=\left(-\frac{4}{7}-1,\ -\frac{5}{7}-2,\ \frac{13}{7}-1\right)
	=\left(-\frac{11}{7},\ -\frac{19}{7},\ \frac{6}{7}\right).
	\]

	\textbf{验证.} $F \in L$：$2\left(-\frac{2}{7}\right)-\frac{5}{14}+\frac{13}{14}=0$，
	$-\frac{2}{7}+\frac{5}{14}+\frac{13}{14}=1$；
	且 $\overrightarrow{P_1 P_2}=\left(-\frac{18}{7},-\frac{33}{7},-\frac{1}{7}\right)$ 与 $\vec{s}$ 的内积为
	$-\frac{36}{7}+\frac{33}{7}+\frac{3}{7}=0$，确实垂直。
}


\item
(本题满分 10 分)
已知 $f(x, y)$ 是 $\left\{(x, y) \mid x^2+y^2 \leqslant 1\right\}$ 上的二次连续可微函数，
满足 $\frac{\partial^2 f}{\partial x^2}+\frac{\partial^2 f}{\partial y^2}=1$，计算
\[
I=\iint\limits_{x^2+y^2 \leqslant 1}\left[\frac{x}{\sqrt{x^2+y^2}} \frac{\partial f}{\partial x}
+\frac{y}{\sqrt{x^2+y^2}} \frac{\partial f}{\partial y}\right] d x d y .
\]

\banswer{
	记 $r=\sqrt{x^2+y^2}$。注意 $\left(\frac{x}{r}, \frac{y}{r}\right)$ 是径向单位向量，故被积函数恰为
	径向导数
	\[
	\frac{x}{r} f_x+\frac{y}{r} f_y=\frac{\partial f}{\partial r} .
	\]
	用极坐标，$I=\displaystyle\int_0^{2 \pi} \int_0^1 f_r(r, \theta)\, r\, d r\, d \theta$。

	极坐标下的 Laplace 算子给出
	\[
	f_{r r}+\frac{1}{r} f_r+\frac{1}{r^2} f_{\theta \theta}=1 .
	\]
	两端乘 $r$，并注意 $r f_{r r}+f_r=\left(r f_r\right)_r$，得
	\[
	\left(r f_r\right)_r+\frac{1}{r} f_{\theta \theta}=r .
	\]
	对 $\theta$ 在 $[0,2 \pi]$ 上积分，由 $f_\theta$ 的周期性 $\int_0^{2 \pi} f_{\theta \theta} d \theta=0$。
	记 $m(r)=\int_0^{2 \pi} f_r(r, \theta) d \theta$，则
	\[
	\frac{d}{d r}[r\, m(r)]=2 \pi r \Longrightarrow r\, m(r)=\pi r^2+C .
	\]
	因 $f \in C^2$，$f_r$ 有界，故 $r \to 0^{+}$ 时 $r\, m(r) \rightarrow 0$，得 $C=0$，即 $m(r)=\pi r$。于是
	\[
	I=\int_0^1 m(r)\, r\, d r=\pi \int_0^1 r^2 d r=\frac{\pi}{3} .
	\]

	\par\vspace{1mm}\noindent
	\textbf{注：}{\kaishu
	可用特例检验：取 $f=\frac{x^2+y^2}{4}$（满足 $\Delta f=1$），则 $f_r=\frac{r}{2}$，
	$I=2 \pi \int_0^1 \frac{r}{2} \cdot r\, d r=\frac{\pi}{3}$；
	再取 $f=\frac{x^2}{2}$，则 $I=\iint \frac{x^2}{r} d x d y=\int_0^{2 \pi} \cos ^2 \theta d \theta \int_0^1 r^2 d r=\frac{\pi}{3}$，
	与一般结论一致。
	}\par
}


\end{enumerate}
